\documentclass[../DoAn.tex]{subfiles}
\begin{document}

\begin{center}
    \Large{\textbf{TÓM TẮT NỘI DUNG ĐỒ ÁN}}\\
\end{center}
\vspace{1cm}
% Sinh viên viết tóm tắt ĐATN của mình trong mục này, với 200 đến 350 từ. Theo trình tự, các nội dung tóm tắt cần có: (i) Giới thiệu vấn đề (tại sao có vấn đề đó, hiện tại được giải quyết chưa, có những hướng tiếp cận nào, các hướng này giải quyết như thế nào, hạn chế là gì), (ii) Hướng tiếp cận sinh viên lựa chọn là gì, vì sao chọn hướng đó, (iii) Tổng quan giải pháp của sinh viên theo hướng tiếp cận đã chọn, và (iv) Đóng góp chính của ĐATN là gì, kết quả đạt được sau cùng là gì. Sinh viên cần viết thành đoạn văn, không được viết ý hoặc gạch đầu dòng.

Nhu cầu quan sát và phân tích mô phỏng giao thông ngày càng tăng, tuy nhiên giao diện 2D truyền thống của SUMO còn hạn chế về khả năng trực quan hóa không gian và theo dõi đồng thời nhiều đối tượng. Một số giải pháp hiển thị 3D đã tồn tại, nhưng thường phụ thuộc vào bản đồ tĩnh, khó tái sử dụng cho nhiều kịch bản và bị giới hạn về số lượng, loại đối tượng giao thông có thể thể hiện. Trong đồ án này, em lựa chọn hướng tiếp cận chuyển đổi dữ liệu mô phỏng từ SUMO sang môi trường hiển thị 3D bằng cách đọc các tệp XML (mạng đường, tuyến, phương tiện, tín hiệu) và dựng cảnh trong Unity3D. Cách tiếp cận này được chọn vì tận dụng được năng lực mô phỏng mạnh của SUMO, đồng thời khai thác pipeline dựng hình và tương tác của Unity3D để tạo một hệ thống hiển thị linh hoạt, có thể mở rộng theo dữ liệu đầu vào. Trên cơ sở đó, đồ án xây dựng quy trình gồm phân tích cấu trúc dữ liệu xuất từ SUMO, thiết kế mô-đun đọc và trích xuất thông tin hình học/topology của mạng, ánh xạ sang hệ tọa độ trong Unity, tạo các đối tượng 3D tương ứng (đường, phương tiện, đèn giao thông) và cập nhật trạng thái theo thời gian mô phỏng để kết xuất chuyển động. Đóng góp chính của đồ án là đề xuất và hiện thực hóa cơ chế dựng bản đồ và đối tượng giao thông dựa trên dữ liệu XML, giúp tự động hóa quá trình tạo cảnh 3D cho nhiều kịch bản; đồng thời xây dựng phiên bản thử nghiệm có thể hiển thị đường giao thông và phương tiện, hỗ trợ quan sát trực quan các tình huống mô phỏng. Kết quả thực nghiệm ban đầu cho thấy hệ thống tái hiện được chuyển động phương tiện trên bản đồ và đáp ứng nhu cầu quan sát, làm tiền đề cho các cải tiến tiếp theo về hiệu năng và chức năng phân tích.

\begin{flushright}
Sinh viên thực hiện\\
\begin{tabular}{@{}c@{}}
\textit{(Ký và ghi rõ họ tên)}
\end{tabular}
\end{flushright}

\end{document}